% GENERATED FILE. Edit canonical lessons, then run npm run generate:books.
% canonical-source-hash: fnv1a64:32134f8d4408cb20
% canonical-lessons: RU-C376-isklyuchenie, RU-C376-lish-by, RU-C376-razve-chto, RU-C376-khotya-by, RU-C376-na-vsyakiy-sluchay, RU-R376-review-369-376

\chapter{Exceptions}
\label{ch:ru-exceptions}
\hlchaptermodality{376}

\begin{chapteropening}
\textbf{By the end of this chapter:} \emph{I can name an exception, say as long as, except perhaps and at least, and do something in case.}

Complete the last lesson of chapter 376: I can name an exception, say as long as, except perhaps and at least, and do something in case.
\end{chapteropening}

\section[\texorpdfstring{\ru{исключение} (isklyuchéniye)}{исключение (isklyuchéniye)}]{\ru{исключение} (isklyuchéniye) — an exception}

\label{lesson:RU-C376-isklyuchenie}



\begin{quote}
Before the new one: say the Russian for necessity, then the Russian for enough.
\end{quote}

\subsection*{You'll want to know: \ru{исключение}}
\textbf{\ru{исключение}} (\emph{isklyuchéniye}) — \textquotedblleft{}an exception\textquotedblright{}.

\begin{quote}
\ru{Все} \ru{пришли}, \ru{за} \ru{исключением} \ru{одного} \ru{студента}.
\end{quote}

\begin{quote}
\emph{Vse prishlí, za isklyuchéniyem odnogó studénta.}
\end{quote}

\begin{quote}
Everyone came except for one student.
\end{quote}

Neuter. \emph{\ru{За} \ru{исключением}} + genitive — except for. \emph{\ru{Без} \ru{исключения}} — without exception.

\begin{cousinweb}
\textbf{\ru{исключение}} is built on \emph{\ru{исключить}} \textquotedblleft{}to exclude\textquotedblright{}, from \emph{\ru{из}-} \textquotedblleft{}out\textquotedblright{} and \emph{\ru{ключ}} \textquotedblleft{}a key\textquotedblright{}, inherited from Proto-Slavic \emph{*ključĭ}: to lock someone out.
\end{cousinweb}

\begin{grammarlens}[title={What you've built}]
One more word for telling, explaining and describing at B1.
\end{grammarlens}

\subsection*{Your turn}
\begin{itemize}
\raggedright
  \item \emph{Say it:} \emph{isklyuchéniye}
  \item \emph{Say it:} \emph{isklyuchéniye}, once more
  \item \emph{Say it:} a sentence of your own with \emph{isklyuchéniye}
  \item \emph{Recall:} say the Russian for necessity, then the Russian for enough, then say \emph{isklyuchéniye} again
\end{itemize}

\begin{tcolorbox}[breakable,colback=teal!4,colframe=teal!35!black,title={Before you move on}]
What does \ru{исключение} mean? (\textquotedblleft{}An exception\textquotedblright{}.) Say it once more.
\end{tcolorbox}
\section[\texorpdfstring{\ru{лишь} \ru{бы} (lish by)}{лишь бы (lish by)}]{\ru{лишь} \ru{бы} (lish by) — as long as, if only}

\label{lesson:RU-C376-lish-by}



\begin{quote}
Before the new one: say the Russian for enough, then the Russian for an exception.
\end{quote}

\subsection*{You'll want to know: \ru{лишь} \ru{бы}}
\textbf{\ru{лишь} \ru{бы}} (\emph{lish by}) — \textquotedblleft{}as long as, if only\textquotedblright{}.

\begin{quote}
\ru{Я} \ru{согласен} \ru{на} \ru{любую} \ru{работу}, \ru{лишь} \ru{бы} \ru{только} \ru{она} \ru{была} \ru{интересной}.
\end{quote}

\begin{quote}
\emph{Ya soglásen na lyubúyu rabótu, lish by tólko oná bylá interésnoy.}
\end{quote}

\begin{quote}
I will take any job, as long as it is interesting.
\end{quote}

\emph{\ru{Лишь} \ru{бы}} + past tense or infinitive names the one thing that matters: \emph{\ru{лишь} \ru{бы} \ru{успеть}}.

\begin{cousinweb}
\textbf{\ru{лишь} \ru{бы}} is built from \emph{\ru{лишь}} \textquotedblleft{}only\textquotedblright{}, which is usually related to \emph{\ru{лишний}} \textquotedblleft{}extra, too many\textquotedblright{}, and \emph{\ru{бы}}, the particle of the conditional.
\end{cousinweb}

\begin{grammarlens}[title={What you've built}]
One more word for telling, explaining and describing at B1.
\end{grammarlens}

\subsection*{Your turn}
\begin{itemize}
\raggedright
  \item \emph{Say it:} \emph{lish by}
  \item \emph{Say it:} \emph{lish by}, once more
  \item \emph{Say it:} a sentence of your own with \emph{lish by}
  \item \emph{Recall:} say the Russian for enough, then the Russian for an exception, then say \emph{lish by} again
\end{itemize}

\begin{tcolorbox}[breakable,colback=teal!4,colframe=teal!35!black,title={Before you move on}]
What does \ru{лишь} \ru{бы} mean? (\textquotedblleft{}As long as, if only\textquotedblright{}.) Say it once more.
\end{tcolorbox}
\section[\texorpdfstring{\ru{разве} \ru{что} (rázve chto)}{разве что (rázve chto)}]{\ru{разве} \ru{что} (rázve chto) — except perhaps, unless maybe}

\label{lesson:RU-C376-razve-chto}



\begin{quote}
Before the new one: say the Russian for an exception, then the Russian for as long as.
\end{quote}

\subsection*{You'll want to know: \ru{разве} \ru{что}}
\textbf{\ru{разве} \ru{что}} (\emph{rázve chto}) — \textquotedblleft{}except perhaps, unless maybe\textquotedblright{}.

\begin{quote}
\ru{Я} \ru{ничего} \ru{не} \ru{хочу}, \ru{разве} \ru{что} \ru{чаю}.
\end{quote}

\begin{quote}
\emph{Ya nichegó ne khochú, rázve chto cháyu.}
\end{quote}

\begin{quote}
I do not want anything, except perhaps some tea.
\end{quote}

It names a small exception: \emph{\ru{никто} \ru{не} \ru{знает}, \ru{разве} \ru{что} \ru{мама}}.

\begin{cousinweb}
\textbf{\ru{разве} \ru{что}} is built from the particle \emph{\ru{разве}} and \emph{\ru{что}} \textquotedblleft{}what\textquotedblright{}. The origin of \emph{\ru{разве}} is not settled.
\end{cousinweb}

\begin{grammarlens}[title={What you've built}]
One more word for telling, explaining and describing at B1.
\end{grammarlens}

\subsection*{Your turn}
\begin{itemize}
\raggedright
  \item \emph{Say it:} \emph{rázve chto}
  \item \emph{Say it:} \emph{rázve chto}, once more
  \item \emph{Say it:} a sentence of your own with \emph{rázve chto}
  \item \emph{Recall:} say the Russian for an exception, then the Russian for as long as, then say \emph{rázve chto} again
\end{itemize}

\begin{tcolorbox}[breakable,colback=teal!4,colframe=teal!35!black,title={Before you move on}]
What does \ru{разве} \ru{что} mean? (\textquotedblleft{}Except perhaps, unless maybe\textquotedblright{}.) Say it once more.
\end{tcolorbox}
\section[\texorpdfstring{\ru{хотя} \ru{бы} (khotyá by)}{хотя бы (khotyá by)}]{\ru{хотя} \ru{бы} (khotyá by) — at least}

\label{lesson:RU-C376-khotya-by}



\begin{quote}
Before the new one: say the Russian for as long as, then the Russian for except perhaps.
\end{quote}

\subsection*{You'll want to know: \ru{хотя} \ru{бы}}
\textbf{\ru{хотя} \ru{бы}} (\emph{khotyá by}) — \textquotedblleft{}at least\textquotedblright{}.

\begin{quote}
\ru{Он} \ru{мог} \ru{бы} \ru{хотя} \ru{бы} \ru{позвонить}.
\end{quote}

\begin{quote}
\emph{On mog by khotyá by pozvonít.}
\end{quote}

\begin{quote}
He could at least have called.
\end{quote}

It sets a minimum: \emph{\ru{хотя} \ru{бы} \ru{немного}} — at least a little. Not the same as \emph{\ru{хотя}}, although.

\begin{cousinweb}
\textbf{\ru{хотя} \ru{бы}} is built from \emph{\ru{хотя}}, the old gerund of \emph{\ru{хотеть}} \textquotedblleft{}to want\textquotedblright{}, and \emph{\ru{бы}}, the particle of the conditional: wanting it to be so, at the very least.
\end{cousinweb}

\begin{grammarlens}[title={What you've built}]
One more word for telling, explaining and describing at B1.
\end{grammarlens}

\subsection*{Your turn}
\begin{itemize}
\raggedright
  \item \emph{Say it:} \emph{khotyá by}
  \item \emph{Say it:} \emph{khotyá by}, once more
  \item \emph{Say it:} a sentence of your own with \emph{khotyá by}
  \item \emph{Recall:} say the Russian for as long as, then the Russian for except perhaps, then say \emph{khotyá by} again
\end{itemize}

\begin{tcolorbox}[breakable,colback=teal!4,colframe=teal!35!black,title={Before you move on}]
What does \ru{хотя} \ru{бы} mean? (\textquotedblleft{}At least\textquotedblright{}.) Say it once more.
\end{tcolorbox}
\section[\texorpdfstring{\ru{на} \ru{всякий} \ru{случай} (na vsyákiy slúchay)}{на всякий случай (na vsyákiy slúchay)}]{\ru{на} \ru{всякий} \ru{случай} (na vsyákiy slúchay) — in case, to be on the safe side}

\label{lesson:RU-C376-na-vsyakiy-sluchay}



\begin{quote}
Before the new one: say the Russian for except perhaps, then the Russian for at least.
\end{quote}

\subsection*{You'll want to know: \ru{на} \ru{всякий} \ru{случай}}
\textbf{\ru{на} \ru{всякий} \ru{случай}} (\emph{na vsyákiy slúchay}) — \textquotedblleft{}in case, to be on the safe side\textquotedblright{}.

\begin{quote}
\ru{Я} \ru{возьму} \ru{зонтик} \ru{на} \ru{всякий} \ru{случай}.
\end{quote}

\begin{quote}
\emph{Ya vozmú zóntik na vsyákiy slúchay.}
\end{quote}

\begin{quote}
I will take an umbrella, to be on the safe side.
\end{quote}

A fixed phrase; \emph{\ru{всякий}} means any.

\begin{cousinweb}
\textbf{\ru{на} \ru{всякий} \ru{случай}} is \textquotedblleft{}for any case\textquotedblright{}: \emph{\ru{всякий}} \textquotedblleft{}any\textquotedblright{} and \emph{\ru{случай}} \textquotedblleft{}a case\textquotedblright{}, inherited from Proto-Slavic \emph{*sŭlučajĭ} (EtymDB).
\end{cousinweb}

\begin{grammarlens}[title={What you've built}]
One more word for telling, explaining and describing at B1.
\end{grammarlens}

\subsection*{Your turn}
\begin{itemize}
\raggedright
  \item \emph{Say it:} \emph{na vsyákiy slúchay}
  \item \emph{Say it:} \emph{na vsyákiy slúchay}, once more
  \item \emph{Say it:} a sentence of your own with \emph{na vsyákiy slúchay}
  \item \emph{Recall:} say the Russian for except perhaps, then the Russian for at least, then say \emph{na vsyákiy slúchay} again
\end{itemize}

\begin{tcolorbox}[breakable,colback=teal!4,colframe=teal!35!black,title={Before you move on}]
What does \ru{на} \ru{всякий} \ru{случай} mean? (\textquotedblleft{}In case, to be on the safe side\textquotedblright{}.) Say it once more.
\end{tcolorbox}
\section[(dialogue)]{Words from chapters 369–376 — a review}

\label{lesson:RU-R376-review-369-376}



\begin{quote}
Say the Russian for ecology and pollution.
\end{quote}

\subsection*{Your turn}
Say these aloud:

\begin{itemize}
\raggedright
  \item ecology, pollution, to pollute, to recycle, to look after
  \item a condition, on condition, in case of, otherwise, even if
  \item a hypothesis, to suppose, let us suppose, let us say, a supposition
  \item probably, probability, unlikely, a chance, a risk
  \item to take a risk, to fear, to avoid, to threaten, to prevent
  \item to count on, to expect, an expectation, to succeed, to work out
  \item to guarantee, definitely, necessary, necessity, enough
  \item an exception, as long as, except perhaps, at least, in case
\end{itemize}

\begin{tcolorbox}[breakable,colback=teal!4,colframe=teal!35!black,title={Before you move on}]
Ecology? (\textbf{\ru{экология}}.) Pollution? (\textbf{\ru{загрязнение}}.) To pollute? (\textbf{\ru{загрязнять}}.) To recycle? (\textbf{\ru{перерабатывать}}.) To look after? (\textbf{\ru{беречь}}.) A condition? (\textbf{\ru{условие}}.) On condition? (\textbf{\ru{при} \ru{условии}}.) In case of? (\textbf{\ru{в} \ru{случае}}.) Otherwise? (\textbf{\ru{в} \ru{противном} \ru{случае}}.) Even if? (\textbf{\ru{даже} \ru{если}}.) A hypothesis? (\textbf{\ru{гипотеза}}.) To suppose? (\textbf{\ru{предполагать}}.) Let us suppose? (\textbf{\ru{предположим}}.) Let us say? (\textbf{\ru{допустим}}.) A supposition? (\textbf{\ru{предположение}}.) Probably? (\textbf{\ru{вероятно}}.) Probability? (\textbf{\ru{вероятность}}.) Unlikely? (\textbf{\ru{маловероятный}}.) A chance? (\textbf{\ru{шанс}}.) A risk? (\textbf{\ru{риск}}.) To take a risk? (\textbf{\ru{рисковать}}.) To fear? (\textbf{\ru{опасаться}}.) To avoid? (\textbf{\ru{избегать}}.) To threaten? (\textbf{\ru{угрожать}}.) To prevent? (\textbf{\ru{предотвращать}}.) To count on? (\textbf{\ru{рассчитывать}}.) To expect? (\textbf{\ru{ожидать}}.) An expectation? (\textbf{\ru{ожидание}}.) To succeed? (\textbf{\ru{удаваться}}.) To work out? (\textbf{\ru{получаться}}.) To guarantee? (\textbf{\ru{гарантировать}}.) Definitely? (\textbf{\ru{обязательно}}.) Necessary? (\textbf{\ru{необходимый}}.) Necessity? (\textbf{\ru{необходимость}}.) Enough? (\textbf{\ru{достаточно}}.) An exception? (\textbf{\ru{исключение}}.) As long as? (\textbf{\ru{лишь} \ru{бы}}.) Except perhaps? (\textbf{\ru{разве} \ru{что}}.) At least? (\textbf{\ru{хотя} \ru{бы}}.) In case? (\textbf{\ru{на} \ru{всякий} \ru{случай}}.)
\end{tcolorbox}
